% vim: set foldmethod=marker foldlevel=0:

\documentclass[a4paper]{article}
\usepackage[UKenglish]{babel}

% \usepackage[hidelinks]{hyperref}

\usepackage{preamble}

% \usepackage{graphicx}
% \graphicspath{ {./imgs/} }

\renewcommand{\thesubsection}{Q\arabic{section}~(\roman{subsection})}
\renewcommand{\thesubsubsection}{(\alph{subsubsection})}

\fancyhead[L]{MA266 Assignment 4}
\title{MA266 Multilinear Algebra, Assignment 4}
\colorlet{questionbodycolor}{magenta!50}

\begin{document}

\maketitle

\setlength{\parindent}{0em}
\setlength{\parskip}{1em}

% {{{ Q14
\question{14}

\begin{questionbody}
Let $\tau \in S_n$. Recall that the permutation matrix associated to $\tau$, written $P_\tau$ is defined as \[
{(P_\tau)}_{i, j} := \begin{cases}
1_{\bb F} & \text{if } j = \tau(i), \\
0_{\bb F} & \text{otherwise}.
\end{cases}
\] Set $L := L(t) := t I_n - P_\tau$ and write $L_{i, j}$ for the $(i, j)$th entry of $L$. Suppose that $\tau$ is the $n$-cycle $\tau = (1, 2, \dotsc, n)$.
\begin{enumerate}[(i)]
\item By using the definition of the determinant from the notes and one of the previous problems on this sheet to compute $\det(L)$, show that $c_{P_\tau}(t) = t^n - 1_{\bb F}$.

\item Suppose that the field $\bb F$ contains an element $\mu \in \bb F$ satisfying $\mu^n = 1_{\bb F}$ and $\mu^m \ne 1_{\bb F}$ for all $1 \le m < n$. This is a primitive $n$th root of unity.
\begin{enumerate}[(a)]
\item Prove that $P_\tau$ is diagonalisable.

\item Determine the minimal polynomial $\mu_{P_\tau}(t)$ of $P_\tau$.
\end{enumerate}
\end{enumerate}
\end{questionbody}

\subsection{~} % 14.i

Answer

\subsection{~} % 14.ii

\subsubsection{~} % 14.ii.a

Answer

\subsubsection{~} % 14.ii.b

Answer

% }}}

% {{{ Q15
\newquestion{15}

\begin{questionbody}
Let $\bb F := \C$. Consider the matrix \[
M := \begin{pmatrix}
7 & 5 & 0 & 0 & 0 \\
-1 & 1 & 0 & 0 & 0 \\
0 & 0 & 0 & 1 & 0 \\
0 & 0 & 0 & 0 & 1 \\
0 & 0 & 1 & 0 & 0
\end{pmatrix} \in M_n(\C).
\]
\begin{enumerate}[(i)]
\item Determine the characteristic polynomial of $M$.

\item Determine the minimal polynomial of $M$.
\end{enumerate}
\end{questionbody}

\subsection{~} % 15.i

Answer

\subsection{~} % 15.ii

Answer

% }}}

\end{document}
